% Minimal preamble for Open Logic Project sources (pinned commit 1e960beff9ed7835bf3e3f1335e21af3439cd107).
%
% OLP's own master preamble cannot run on the project host: TeX Live 2025 there lacks
% prooftree.sty, ntheorem.sty, extarrows.sty, mathdots.sty, isomath.sty and subfiles.sty.
% Following the project's own rule ("if the author preamble needs packages the toolchain lacks,
% do not invent macros -- use a minimal working preamble"), this file declares, verbatim in
% spirit and numbering behaviour, exactly the theorem-like environments OLP defines in
% sty/open-logic-book-envs.sty (thm/lem/prop/cor/defn/ex/prob/rem/note/case/axiom/conv),
% plus the shorthand accents OLP uses (`!!a`, `!!b`, `!!{...}`).
%
% Author text is untouched; only the packaging is ours.
\documentclass[10pt,a4paper]{article}

\usepackage{amsmath}
\usepackage{amssymb}
\usepackage{amsthm}
\usepackage{graphicx}
\usepackage{hyperref}
\usepackage{enumitem}
\usepackage[framemethod=default]{mdframed}

% ---- OLP shorthand accents (from sty/open-logic.sty) ----
% NB: \! is taken by LaTeX (negative thin space), so we only define the !!a / !!b families
% and the \olup helper used by them.
\newcommand*{\olup}[1]{\MakeUppercase{#1}}
\def\!!!a#1{\olup{#1}}
\def\!a#1{\olup{#1}}
\def\!b#1{\olup{#1}}

% ---- theorem-like environments, same family numbering as OLP ----
\theoremstyle{plain}
\newtheorem{thm}{Theorem}[section]
\newtheorem{lem}[thm]{Lemma}
\newtheorem{prop}[thm]{Proposition}
\newtheorem{cor}[thm]{Corollary}
\newtheorem{axiom}{Axiom}

\theoremstyle{definition}
\newtheorem{defn}[thm]{Definition}
\newtheorem{ex}[thm]{Example}
\newtheorem{prob}{Problem}[section]
\newtheorem{probd}{Problem}
\newtheorem{conv}[thm]{Convention}

\theoremstyle{remark}
\newtheorem{rem}[thm]{Remark}
\newtheorem{note}[thm]{Note}
\newtheorem{case}[thm]{Case}

% ---- environments OLP defines as wrappers ----
\newenvironment{probtag}[1]{\begin{prob}}{\end{prob}}
\newenvironment{explain}{}{}
\newenvironment{explain*}{}{}
\newenvironment{editorial}{}{}
\newenvironment{digress}{}{}
\newenvironment{digress*}{}{}
\newenvironment{history}{}{}
\newenvironment{reference}{}{}
\newcommand*{\olfileid}[3]{}
\newcommand*{\olsection}[1]{\section*{#1}}
\newcommand*{\olchapter}[3]{}
\newcommand*{\olpart}[2]{}
\newcommand*{\olimport}[1]{}
\newcommand*{\olinclude}[1]{}
\newcommand*{\olpath}[1]{}
\newcommand*{\ollangid}[1]{}
\newcommand*{\problemtag}[1]{}
\newcommand*{\problemauthor}[1]{}
\newcommand*{\ollabel}[1]{\label{#1}}
% \olref takes 0-3 OPTIONAL arguments plus one mandatory label:
%     \olref{lab}  |  \olref[from]{lab}  |  \olref[from][to]{lab}  |  \olref[a][b][c]{lab}
% A single-argument definition silently swallows "[sth]" and emits a broken reference
% (measured on the host: the text came out as "??iso]x"). The chain below consumes the
% brackets and cross-references only the mandatory argument.
\makeatletter
\newcommand{\olref}{\@ifnextchar[{\olref@i}{\olref@final}}
\def\olref@i[#1]{\@ifnextchar[{\olref@ii}{\olref@final}}
\def\olref@ii[#1]{\@ifnextchar[{\olref@iii}{\olref@final}}
\def\olref@iii[#1]{\olref@final}
\def\olref@final#1{\ref{#1}}
\makeatother
\newcommand*{\olnewpage}{\clearpage}

% ===================== OLP compatibility layer (1) =====================

\makeatletter

% ---------------- proof trees (OLP's prooftree is a bussproofs derivative) ----------------
\providecommand{\AxiomC}[1]{}
\providecommand{\UnaryInfC}[1]{}
\providecommand{\BinaryInfC}[1]{}
\providecommand{\TrinaryInfC}[1]{}
\providecommand{\QuaternaryInfC}[1]{}
\providecommand{\QuinaryInfC}[1]{}
\providecommand{\RightLabel}[1]{}
\providecommand{\LeftLabel}[1]{}
\providecommand{\doubleLine}{}
\providecommand{\noLine}{}
\providecommand{\DisplayProof}{}
\providecommand{\fCenter}{}
\newenvironment{prooftree}{}{}
\newenvironment{ltprooftree}{}{}
\newenvironment{oltableau}{}{}

% ---------------- logical notation ----------------
\providecommand{\lif}{\mathbin{\rightarrow}}
\providecommand{\lforall}{\mathord{\forall}}
\providecommand{\lexists}{\mathord{\exists}}
\providecommand{\ltrue}{\top}
\providecommand{\lfalse}{\bot}
\providecommand{\Ltrue}{\top}
\providecommand{\Lfalse}{\bot}
\providecommand{\Sequent}{\mathrel{\Rightarrow}}
\providecommand{\Proves}{\mathrel{\vdash}}
\providecommand{\Entails}{\mathrel{\models}}
\providecommand{\Subst}[3]{#1[#2/#3]}
\providecommand{\sFmla}[2]{#1(#2)}
\providecommand{\pFmla}[2]{#1(#2)}
\providecommand{\formula}{\relax}
\providecommand{\True}{\top}
\providecommand{\False}{\bot}
\providecommand{\Args}[1]{#1}
\providecommand{\Ts}{\ensuremath{\top}}
\providecommand{\Fi}{\ensuremath{\bot}}
\providecommand{\lfalsum}{\bot}

% ---------------- objects, structures, metatheory ----------------
\providecommand{\Obj}{}
\providecommand{\Struct}{\ensuremath{\mathfrak{A}}}
\providecommand{\Domain}{\ensuremath{|\mathfrak{A}|}}
\providecommand{\Assign}{\ensuremath{s}}
\providecommand{\Sat}[2]{\ensuremath{#1 \models #2}}
\providecommand{\mSat}[3]{\ensuremath{#1 \models #2[#3]}}
\providecommand{\mModel}[2]{\ensuremath{#1 \models #2}}
\providecommand{\Lang}{\ensuremath{\mathcal{L}}}
\providecommand{\Th}{\ensuremath{\mathrm{Th}}}
\providecommand{\Log}{\ensuremath{\mathrm{L}}}
\providecommand{\tuple}[1]{\langle #1 \rangle}
\providecommand{\num}[1]{\underline{#1}}
\providecommand{\cfind}[1]{\varphi_{#1}}
\providecommand{\Nat}{\ensuremath{\mathbb{N}}}
\providecommand{\Mult}{\ensuremath{\mathrm{Mult}}}
\providecommand{\Add}{\ensuremath{\mathrm{Add}}}
\providecommand{\Exp}{\ensuremath{\mathrm{Exp}}}
\providecommand{\Zero}{\ensuremath{0}}
\providecommand{\Succ}{\ensuremath{\mathrm{S}}}
\providecommand{\lambd}{\ensuremath{\lambda}}
\providecommand{\bredpar}{\mathrel{\rightarrow_\beta}}
\providecommand{\bcd}[1]{\ensuremath{\ulcorner #1 \urcorner}}
\providecommand{\Set}[1]{\ensuremath{\{#1\}}}
\providecommand{\indcase}[1]{\par\noindent\emph{Case #1}.\quad}
\providecommand{\indfrm}[1]{\par\noindent\emph{Form #1}.\quad}
\providecommand{\iftag}[1]{#1}
\providecommand{\tagitem}[1]{\item[#1]}
\providecommand{\Elim}[1]{}
\providecommand{\Intro}[1]{}
\providecommand{\LeftR}[1]{}
\providecommand{\RightR}[1]{}
\providecommand{\TAss}[1]{}
\providecommand{\TRule}[1]{}

\makeatother
\providecommand{\olpath}{.}

% ===================== OLP compatibility layer (2) =====================
% --- second compatibility block: notation and shorthands used by the frozen OLP units ---
\providecommand{\CutCS}{\mathrm{Cut}}
\providecommand{\LeftSubproofLabel}{\relax}
\providecommand{\RightSubproofLabel}{\relax}
\providecommand{\maeh}{\relax}
\providecommand{\cheight}{\relax}
\providecommand{\ande}{\mathrm{and}_E}
\providecommand{\andi}{\mathrm{and}_I}
\providecommand{\ore}{\mathrm{or}_E}
\providecommand{\ori}{\mathrm{or}_I}
\providecommand{\typeof}{\mathord{:}}
\providecommand{\pValue}[1]{#1}
\providecommand{\Sent}{\mathrm{Sent}}
\providecommand{\cardplus}{\mathbin{+}}
\providecommand{\cardtimes}{\mathbin{\cdot}}
\providecommand{\cardexpo}{\mathbin{\hat{\ }}}
\providecommand{\cardnless}{\not<}
\providecommand{\cardle}{\leq}
\providecommand{\funfromto}[3]{#1\colon #2 \to #3}
\providecommand{\funimage}[1]{\operatorname{im}(#1)}
\providecommand{\fdefined}{\mathord{\downarrow}}
\providecommand{\canonord}{\operatorname{ord}}
\providecommand{\ordtimes}{\mathbin{\cdot}}
\providecommand{\ordeq}[2]{#1 = #2}
\providecommand{\ordneq}[2]{#1 \neq #2}
\providecommand{\ordtype}[1]{\operatorname{ord}(#1)}
\providecommand{\ordsucc}{\mathord{+1}}
\providecommand{\ordunion}{\cup}
\providecommand{\onesphere}{S^1}
\providecommand{\rotationsgroup}{\mathrm{SO}(2)}
\providecommand{\equivrep}[1]{[#1]}
\providecommand{\rlexless}{\prec}
\providecommand{\ZFC}{\mathrm{ZFC}}
\providecommand{\ZF}{\mathrm{ZF}}
\providecommand{\Rat}{\mathbb{Q}}
\providecommand{\Real}{\mathbb{R}}
\providecommand{\PosInt}{\mathbb{Z}^{+}}
\providecommand{\Id}{\mathrm{id}}
\providecommand{\Bin}{\mathrm{Bin}}
\providecommand{\pto}{\rightharpoonup}
\providecommand{\dom}{\operatorname{dom}}
\providecommand{\ran}{\operatorname{ran}}
\providecommand{\len}{\operatorname{len}}
\providecommand{\concat}{\mathbin{{}^\frown}}
\providecommand{\card}[1]{\lvert #1 \rvert}
\providecommand{\Pow}{\mathcal{P}}
\providecommand{\Setabs}[2]{\{\, #1 \mid #2 \,\}}
\providecommand{\comp}{\circ}
\providecommand{\Char}{\chi}
\providecommand{\Prf}{\operatorname{Prf}}
\providecommand{\Prov}{\operatorname{Prov}}
\providecommand{\Proj}{\operatorname{proj}}
\providecommand{\Gn}{\operatorname{gn}}
\providecommand{\gn}{\operatorname{gn}}
\providecommand{\Value}[1]{\operatorname{val}(#1)}
\providecommand{\Deduce}{\mathrm{D}}
\providecommand{\Dual}{\mathrm{d}}
\providecommand{\Ax}{\mathrm{Ax}}
\providecommand{\pSat}[3]{#1 \models #2\,[#3]}
\providecommand{\pAssign}{\sigma}
\providecommand{\bforall}{\forall}
\providecommand{\bexists}{\exists}
\providecommand{\liff}{\mathbin{\leftrightarrow}}
\providecommand{\FV}{\operatorname{FV}}
\providecommand{\UnaryInf}{\ensuremath{\mid}}
\providecommand{\BinaryInf}{\ensuremath{\mid}}
\providecommand{\depth}{\operatorname{depth}}
\providecommand{\oliflabeldef}[3]{#1}
\providecommand{\TMblank}{\sqcup}
\providecommand{\TMstroke}{\mid}
\providecommand{\TMstay}{\mathrm{S}}
\providecommand{\TMright}{\mathrm{R}}
\providecommand{\TMleft}{\mathrm{L}}
\providecommand{\TMendtape}{\triangleright}
\providecommand{\umin}{\mu}
\providecommand{\Sum}{\sum}
\providecommand{\citep}[1]{\cite{#1}}
\providecommand{\citet}[1]{\cite{#1}}
\providecommand{\citeauthor}[1]{\cite{#1}}
\providecommand{\citeyear}[1]{\cite{#1}}

% ===================== OLP compatibility layer (3) =====================
% --- third compatibility block: remaining notation seen in the frozen OLP units ---
\providecommand{\iso}{\cong}
\providecommand{\nszero}{\mathrm{0}}
\providecommand{\nssucc}{\mathrm{S}}
\providecommand{\nsplus}{\mathbin{+}}
\providecommand{\nssless}{<}
\providecommand{\nsless}{<}
\providecommand{\PL}{\mathrm{PL}}
\providecommand{\Nec}{\mathord{\Box}}
\providecommand{\RK}{\mathrm{RK}}
\providecommand{\ST}{\mathrm{ST}}
\providecommand{\SSubst}[3]{#1[#2/#3]}
\providecommand{\subst}[3]{#1[#2/#3]}
\providecommand{\tagfalse}{\mathrm{F}}
\providecommand{\cutr}{\mathrm{Cut}}
\providecommand{\cardneq}{\neq}
\providecommand{\issucc}[1]{\operatorname{succ}(#1)}
\providecommand{\tup}[1]{\langle #1 \rangle}
\providecommand{\Var}{\mathrm{Var}}
\providecommand{\Bew}{\operatorname{Bew}}
\providecommand{\Con}{\operatorname{Con}}
\providecommand{\Cons}{\operatorname{Cons}}
\providecommand{\pDeduce}{\mathrm{D}}
\providecommand{\psat}{\models}
\providecommand{\den}[1]{\llbracket #1 \rrbracket}
\providecommand{\Sem}[2]{\llbracket #1 \rrbracket_{#2}}
\providecommand{\DomainOf}[1]{\lvert #1 \rvert}
\providecommand{\Form}{\mathrm{Form}}
\providecommand{\Terms}{\mathrm{Term}}
\providecommand{\Atom}{\mathrm{Atom}}
\providecommand{\FVars}[1]{\operatorname{FV}(#1)}
\providecommand{\fv}[1]{\operatorname{fv}(#1)}
\providecommand{\tuplep}[1]{\langle #1 \rangle}
\providecommand{\restrict}[2]{#1 \upharpoonright #2}
\providecommand{\sto}{\rightarrow}
\providecommand{\FmlaSet}[1]{\mathrm{Fmla}_{#1}}
\providecommand{\Thm}{\mathrm{Thm}}
\providecommand{\Eval}{\operatorname{val}}
\providecommand{\blank}{\sqcup}
\providecommand{\blankTape}{\sqcup}
\providecommand{\Consis}{\operatorname{Con}}
\providecommand{\drs}{\mathrm{d}}
\providecommand{\VRule}{\mid}
\providecommand{\Bet}{\mathrm{B}}
\providecommand{\Seg}[1]{\mathrm{Seg}_{#1}}
\providecommand{\Enums}{\mathrm{En}}
\providecommand{\Z}{\mathbb{Z}}
\providecommand{\N}{\mathbb{N}}
\providecommand{\abs}[1]{\lvert #1 \rvert}
\providecommand{\floor}[1]{\lfloor #1 \rfloor}
\providecommand{\ceil}[1]{\lceil #1 \rceil}

% ===================== OLP compatibility layer (4) =====================
% --- fourth compatibility block: harvested from the per-candidate compile logs ---
\providecommand{\beredpar}{\mathrel{\rightarrow_\beta}}
\makeatletter
\def\bered{\@ifnextchar[{\bered@i}{\mathrel{\rightarrow_\beta}}}
\def\bered@i[#1]{\mathrel{\rightarrow_\beta}}
\def\bred{\@ifnextchar[{\bred@i}{\mathrel{\rightarrow_\beta}}}
\def\bred@i[#1]{\mathrel{\rightarrow_\beta}}
\makeatother
\providecommand{\beredone}{\mathrel{=_\beta}}
\providecommand{\bredone}{\mathrel{=_\beta}}
\providecommand{\xredone}{\mathrel{=_\beta}}
\providecommand{\aconvone}{\mathrel{=_\alpha}}
\providecommand{\cardless}{<}
\providecommand{\tagrefs}[1]{#1}
\providecommand{\ordplus}{\mathbin{+}}
\providecommand{\funrestrictionto}[2]{#1 \upharpoonright #2}
\providecommand{\fundefined}{\mathord{\uparrow}}
\providecommand{\Trm}{\mathrm{Trm}}
\providecommand{\OProv}{\operatorname{Prov}}
\providecommand{\OPrf}{\operatorname{Prf}}
\providecommand{\aeq}{\mathrel{=}}
\providecommand{\mClass}[1]{[#1]}
\providecommand{\equivclass}[1]{[#1]}
\providecommand{\pheight}{\operatorname{ht}}
\providecommand{\phy}{\operatorname{ht}}
\providecommand{\Cut}{\mathrm{Cut}}
\providecommand{\QuantRank}{\operatorname{qr}}
\providecommand{\elemequiv}{\mathrel{\equiv}}
\providecommand{\PIso}{\mathrm{PIso}}
\providecommand{\Theory}{\mathrm{Th}}
\providecommand{\supstrict}{\subsetneq}
\providecommand{\trcl}{\operatorname{trcl}}
\providecommand{\cref}[1]{\ref{#1}}
\providecommand{\Discharge}{\mathrm{Dis}}
\providecommand{\Weakening}{\mathrm{Wk}}
\providecommand{\ZFminus}{\mathrm{ZF}^{-}}
\makeatletter
\renewcommand{\Dual}{\@ifnextchar[{\Dual@i}{\ifmmode\mathrm{d}\else$\mathrm{d}$\fi}}
\def\Dual@i[#1]{\ifmmode\mathrm{d}[#1]\else$\mathrm{d}[#1]$\fi}
\makeatother
\providecommand{\RightR}{\mathrm{R}}
\providecommand{\LeftR}{\mathrm{L}}
\providecommand{\Intro}{\mathrm{I}}
\providecommand{\Axiom}{\mathrm{Ax}}
\providecommand{\formula}[1]{#1}
\providecommand{\stared}{\ast}
\providecommand{\oset}[2]{\{\, #1 \mid #2 \,\}}
\providecommand{\setof}[2]{\{\, #1 \mid #2 \,\}}
\providecommand{\pow}[1]{\mathcal{P}(#1)}
\providecommand{\ordpair}[2]{\langle #1, #2 \rangle}
\providecommand{\ordtriple}[3]{\langle #1, #2, #3 \rangle}
\providecommand{\union}{\cup}
\providecommand{\inter}{\cap}
\providecommand{\seq}[1]{\langle #1 \rangle}
\providecommand{\SuchThat}{\mid}
\providecommand{\where}{\mid}
\providecommand{\TVar}{\mathrm{TVar}}
\providecommand{\Bool}{\mathrm{Bool}}
\providecommand{\Val}{\operatorname{Val}}
\providecommand{\TL}[1]{\mathrm{TL}}
\providecommand{\xstared}{\ast}

% ===================== OLP compatibility layer (5) =====================
% --- fifth compatibility block: rescues the quarantined OLP candidates ---
\providecommand{\ident}{\operatorname{id}}
\providecommand{\setrank}{\operatorname{rank}}
\providecommand{\cardfont}{\mathfrak{c}}
\providecommand{\Complement}{\overline}
\providecommand{\QR}{\mathrm{QR}}
\providecommand{\becd}{\ensuremath{\ulcorner}}
\providecommand{\disjointsum}{\mathbin{\oplus}}
\providecommand{\cardeq}[2]{#1 = #2}
\providecommand{\xred}{\mathrel{\rightarrow}}
% --- sixth compatibility block: names reported by real compile logs ---
\providecommand{\ORefut}{\operatorname{Refut}}
\providecommand{\ORProv}{\operatorname{Prov}}
\providecommand{\cutrank}{\operatorname{cr}}
\providecommand{\Refut}{\operatorname{Refut}}
\providecommand{\DischargeRule}{\mathrm{Dis}}
\providecommand{\Contraction}{\mathrm{C}}
\providecommand{\bottomAlignProof}{\relax}
\providecommand{\FalseCl}{\bot}
\providecommand{\red}{\mathrel{\rightarrow}}
\providecommand{\varAssign}{\sigma}
\providecommand{\defiff}{\mathrel{\leftrightarrow}}
\providecommand{\FalseInt}{\bot}
\providecommand{\fn}{\operatorname{fn}}
\providecommand{\dotted}{\cdot}

% --- tag-macro semantic fix (2026-10-13, final): \iftag CHOOSES a branch and the OLP proof body is
% always the SECOND group (\iftag{tag}{body}, {tag}{body}{c}, {tag}{body}{c}{d}); a definition that
% prints the tag literal drops the proof body. Bodies contain blank lines, hence \long.
\makeatletter
\providecommand{\iftag}{}
\long\def\iftag#1#2{\@ifnextchar\bgroup{\iftag@three{#2}}{\iftag@stop{#2}}}
\long\def\iftag@three#1#2{\iftag@stop{#1}\@ifnextchar\bgroup{\iftag@eat}{\relax}}
\long\def\iftag@eat#1{}
\long\def\iftag@stop#1{#1}
\long\def\oliflabeldef#1#2{\@ifnextchar\bgroup{\@gobble}{\iftag@stop{#2}}}
\providecommand{\tagrefs}[1]{}
\providecommand{\settag}[1]{}
\makeatother
\begin{document}

% ==== OLP file: content/model-theory/basics/basics.tex ====
% Part: model-theory
% Chapter: basics



\olchapter{mod}{bas}{Basics of Model Theory}

\olimport{reducts-and-expansions}

\olimport{substructures}

\olimport{overspill}

\olimport{isomorphism}

\olimport{theory-of-m}

\olimport{partial-iso}

\olimport{dlo}

%\olimport{nonstandard-arithmetic}

\OLEndChapterHook


% ==== OLP file: content/model-theory/basics/dlo.tex ====
% Part: first-order-logic
% Chapter: model-theory
% Section: partial-iso



\olfileid{mod}{bas}{dlo}
\section{Dense Linear Orders}

\begin{defn}
  A \emph{dense linear ordering without endpoints} is \olup{structure}
  $\Struct{M}$ for the \olup{language} containing a single 2-place
  \olup{predicate}~$<$ satisfying the following sentences:
  \begin{enumerate}
  \item $\lforall[x][\lnot x < x]$;
  \item $\lforall[x][\lforall[y][\lforall[z][(x < y \lif (y < z \lif x
    <z ))]]]$;
  \item $\lforall[x][\lforall[y][(x< y \lor \eq[x][y] \lor y < x)]]$;
  \item $\lforall[x][\lexists[y][x < y]]$;
  \item $\lforall[x][\lexists[y][y < x]]$;
  \item $\lforall[x][\lforall[y][(x < y \lif \lexists[z][(x < z \land
        z < y))]]]$.
 \end{enumerate}
\end{defn}

\begin{thm}\ollabel{thm:cantorQ}
  Any two \olup{enumerable} dense linear orderings without
  endpoints are isomorphic.
\end{thm}

\begin{proof}
  Let $\Struct{M_1}$ and $\Struct{M_2}$ be \olup{enumerable} dense linear
  orderings without endpoints, with ${<_1} = \Assign{<}{M_1}$ and ${<_2} =
  \Assign{<}{M_2}$, and let $\PIso{I}$ be the set of all partial
  isomorphisms between them. $\PIso{I}$ is not empty since at least
  $\emptyset \in \PIso{I}$. We show that $\PIso{I}$ satisfies the
  Back-and-Forth property.  Then $\Struct{M_1} \iso[p] \Struct{M_2}$,
  and the theorem follows by \olref[pis]{thm:p-isom1}.

  To show $\PIso{I}$ satisfies the Forth property, let $p \in
  \PIso{I}$ and let $p(a_i) = b_i$ for $i = 1$, \dots,~$n$, and
  without loss of generality suppose $a_1 <_1 a_2 <_1 \cdots <_1
  a_n$. Given $a \in \Domain{M_1}$, find $b \in \Domain{M_2}$ as
  follows:
  \begin{enumerate}
  \item if $a <_1 a_1$ let $b \in \Domain{M_2}$ be such that $b <_2
    b_1$;
  \item if $a_n <_1 a$ let $b \in \Domain{M_2}$ be such that $b_n <_2 b$;
 \item if $a_i <_1 a <_1 a_{i+1}$ for some $i$, then let $b \in
   \Domain{M_2}$ be such that $b_i <_2 b <_2 b_{i+1}$.
  \end{enumerate}
  It is always possible to find a $b$ with the desired property since
  $\Struct{M_2}$ is a dense linear ordering without endpoints. Define
  $q = p \cup \{ \langle a, b \rangle \}$ so that $q \in \PIso{I}$ is
  the desired extension of $p$. This establishes the Forth
  property. The Back property is similar. So $\Struct{M_1} \iso[p]
  \Struct{M_2}$; by \olref[pis]{thm:p-isom1}, $\Struct{M_1} \iso
  \Struct{M_2}$.
\end{proof}

\begin{prob}
  Complete the proof of \olref[mod][bas][dlo]{thm:cantorQ} by
  verifying that $\PIso{I}$ satisfies the Back property.
\end{prob}

\begin{rem}
  Let $\Struct{S}$ be any \olup{enumerable} dense linear ordering without
  endpoints. Then (by \olref{thm:cantorQ}) $\Struct{S} \iso
  \Struct{Q}$, where $\Struct{Q} = (\Rat, <)$ is the \olup{enumerable}
  dense linear ordering having the set $\Rat$ of the rational numbers
  as its domain. Now consider again the \olup{structure}~$\Struct{R} =
  (\Real, <)$ from \olref[thm]{remark:R}. We saw that there is
  \olup{enumerable} \olup{structure}~$\Struct{S}$ such that $\Struct{R}
  \elemequiv \Struct{S}$. But $\Struct{S}$ is \olup{enumerable} dense
  linear ordering without endpoints, and so it is isomorphic (and
  hence elementarily equivalent) to the \olup{structure}~$\Struct{Q}$. By
  transitivity of elementary equivalence, $\Struct{R} \elemequiv
  \Struct{Q}$. (We could have shown this directly by establishing
  $\Struct{R} \iso[p] \Struct{Q}$ by the same back-and-forth
  argument.)
\end{rem}

% ==== OLP file: content/model-theory/basics/isomorphism.tex ====
% Part: first-order-logic
% Chapter: model-theory
% Section: isomorphism



\olfileid{mod}{bas}{iso}
\olsection{Isomorphic Structures}

First-order \olup{structure}s can be alike in one of two ways. One way in
which they can be alike is that they make the same \olup{sentence}s
true. We call such \olup{structure}s \emph{elementarily equivalent}. But
structures can be very different and still make the same \olup{sentence}s
true---for instance, one can be \olup{enumerable} and the other not.
This is because there are lots of features of \olup{structure} that
cannot be expressed in first-order languages, either because the
language is not rich enough, or because of fundamental limitations of
first-order logic such as the L\"owenheim--Skolem theorem. So another,
stricter, aspect in which \olup{structure}s can be alike is if they are
fundamentally the same, in the sense that they only differ in the
objects that make them up, but not in their structural features. A way
of making this precise is by the notion of an \emph{isomorphism}.

\begin{defn}
\ollabel{defn:elem-equiv} 
Given two \olup{structure}s $\Struct{M}$ and $\Struct M'$ for the same
\olup{language}~$\Lang{L}$, we say that $\Struct{M}$ is \emph{elementarily
  equivalent to} $\Struct M'$, written $\Struct{M} \equiv \Struct M'$,
if and only if for every \olup{sentence}~$A$ of~$\Lang{L}$,
$\Sat{M}{A}$ iff $\Sat{M'}{A}$.
\end{defn}

\begin{defn}
\ollabel{defn:isomorphism} Given two \olup{structure}s $\Struct{M}$ and
$\Struct M'$ for the same \olup{language}~$\Lang L$, we say that
$\Struct{M}$ is \emph{isomorphic to}~$\Struct M'$, written $\Struct{M}
\simeq \Struct M'$, if and only if there is a function $h \colon
\Domain{M} \to \Domain{M'}$ such that:
\begin{enumerate}
\item $h$ is \olup{injective}: if $h(x) =
  h(y)$ then $x = y$; 
\item $h$ is \olup{surjective}: for every $y \in \Domain{M'}$ there
  is $x \in \Domain{M}$ such that $h(x) = y$;
\item \ollabel{defn:iso-const}for every \olup{constant} $c$:
  $h(\Assign{c}{M}) = \Assign{c}{M'}$;
\item \ollabel{defn:iso-pred}for every $n$-place \olup{predicate}~$P$:
  \[
  \tuple{a_1, \dots, a_n}\in \Assign{P}{M} \quad\text{iff}\quad
  \tuple{h(a_1), \dots, h(a_n)} \in \Assign{P}{M'};
  \]
\item \ollabel{defn:iso-func}for every $n$-place \olup{function} $f$:
  \[
  h(\Assign{f}{M}(a_1, \dots, a_n)) =
  \Assign{f}{M'}(h(a_1), \dots, h(a_n)).
  \]
\end{enumerate}
\end{defn}

\begin{thm}
\ollabel{thm:isom}
If $\Struct{M} \iso \Struct M'$ then $\Struct{M} \elemequiv
\Struct{M'}$.
\end{thm}

\begin{proof}
Let $h$ be an isomorphism of $\Struct{M}$ onto $\Struct M'$. For any
assignment~$s$, $h \circ s$ is the composition of $h$ and $s$, i.e.,
the assignment in $\Struct{M'}$ such that $(h \circ s)(x) = h(s(x))$.
By induction on $t$ and $A$ one can prove the stronger claims:
\begin{enumerate}
  \item[a.] $h(\Value{t}{M}[s]) = \Value{t}{M'}[h\circ s]$.
  \item[b.] $\Sat{M}{A}[s]$ iff $\Sat{M'}{A}[h \circ s]$.
\end{enumerate}
The first is proved by induction on the complexity of~$t$.
\begin{enumerate}
\item If $t \ident c$, then $\Value{c}{M}[s] = \Assign{c}{M}$ and
  $\Value{c}{M'}[h \circ s] = \Assign{c}{M'}$. Thus,
  $h(\Value{t}{M}[s]) = h(\Assign{c}{M}) = \Assign{c}{M'}$ (by
  \olref{defn:iso-const} of \olref{defn:isomorphism}) $=
  \Value{t}{M'}[h \circ s]$.
\item If $t \ident x$, then $\Value{x}{M}[s] = s(x)$ and
  $\Value{x}{M'}[h \circ s] = h(s(x))$. Thus, $h(\Value{x}{M}[s]) =
  h(s(x)) = \Value{x}{M'}[h \circ s]$.
\item If $t \ident f(t_1, \dots, t_n)$, then
  \begin{align*}
    \Value{t}{M}[s] & = \Assign{f}{M}(\Value{t_1}{M}[s], \dots, \Value{t_n}{M}[s]) \quad\text{and}\\
  \Value{t}{M'}[h \circ s] & = \Assign{f}{M'}(\Value{t_1}{M'}[h \circ
    s], \dots, \Value{t_n}{M'}[h \circ s]).
  \end{align*}
  The induction hypothesis is that for each $i$, $h(\Value{t_i}{M}[s])
  = \Value{t_i}{M'}[h\circ s]$. So,
  \begin{align}
    h(\Value{t}{M}[s]) 
    & = h(\Assign{f}{M}(\Value{t_1}{M}[s], \dots, \Value{t_n}{M}[s]))\notag\\
    & = \Assign{f}{M'}(h(\Value{t_1}{M}[s]), \dots,
    h(\Value{t_n}{M}[s])) \ollabel{iso-1}\\
    & = \Assign{f}{M'}(\Value{t_1}{M'}[h \circ s], \dots,
    \Value{t_n}{M'}[h \circ s]) \ollabel{iso-2}\\
    & = \Value{t}{M'}[h\circ s] \notag
  \end{align}
  Here, \olref{iso-1} follows by \olref{defn:iso-func} of
  \olref{defn:isomorphism} and \olref{iso-2} by induction hypothesis.
\end{enumerate}
Part (b) is left as an exercise.

If $A$ is a sentence, the assignments~$s$ and $h \circ s$ are
irrelevant, and we have $\Sat{M}{A}$ iff $\Sat{M'}{A}$.
\end{proof}

\begin{prob}
Carry out the proof of (b) of \olref[mod][bas][iso]{thm:isom} in
detail. Make sure to note where each of the five properties
characterizing isomorphisms of \olref[mod][bas][iso]{defn:isomorphism}
is used.
\end{prob}

\begin{defn}
An \emph{automorphism} of a structure $\Struct{M}$ is an isomorphism
of $\Struct{M}$ onto itself.
\end{defn}

\begin{prob}
Show that for any structure $\Struct{M}$, if $X$ is a definable subset
of $\Struct{M}$, and $h$ is an automorphism of $\Struct{M}$, then $X
= \Setabs{h(x)}{x \in X}$ (i.e., $X$ is fixed under $h$).
\end{prob}



% ==== OLP file: content/model-theory/basics/nonstandard-arithmetic.tex ====
% Part: first-order-logic
% Chapter: model-theory
% Section: nonstandard-arithmetic



\olfileid{mod}{bas}{nsa}
\section{Non-standard Models of Arithmetic}

\begin{defn}
Let $\Lang{L_N}$ be the \olup{language} of arithmetic, comprising a
\olup{constant} $\Obj{0}$, a 2-place \olup{predicate} $<$, a 1-place
\olup{function} $\prime$, and 2-place \olup{function}s $+$ and $\times$.
\begin{enumerate}
\item The \emph{standard model} of arithmetic is the \olup{structure}
  $\Struct{N}$ for $\Lang{L_N}$ having $\Nat = \{ 0, 1, 2, \dots\}$
  and interpreting $\Obj{0}$ as $0$, $<$ as the less-than relation
  over~$\Nat$, and $\prime$, $+$ and $\times$ as successor, addition,
  and multiplication over $\Nat$, respectively.
\item \emph{True arithmetic} is the theory $\Theory{N}$.
\end{enumerate}
\end{defn}

When working in $\Lang{L_N}$ we abbreviate each term of
the form $\Obj{0}^{\prime\dots\prime}$, with $n$ applications of the
successor function to $\Obj{0}$, as~$\num{n}$.

\begin{defn}
A \olup{structure}~$\Struct{M}$ for $\Lang{L_N}$ is \emph{standard} if
and only $\Struct{N} \iso \Struct{M}$.
\end{defn}

\begin{thm}\ollabel{thm:non-std}
There are non-standard \olup{enumerable} models of true arithmetic.
\end{thm}

\begin{proof}
Expand $\Lang{L_N}$ by introducing a new \olup{constant}~$c$, and
consider the theory
\[
\Theory{N} \cup \Setabs{\num{n} < c}{n \in \Nat}.
\]
The theory is finitely satisfiable, so by compactness it has a model
$\Struct{M}$, which can be taken to be \olup{enumerable} by the Downward
L\"owenheim--Skolem theorem. Where $\Domain{M}$ is the domain of
$\Struct{M}$, let $\Struct{M}$ interpret the non-logical constants of
$\Lang{L}$ as $\mathbf{z} = \Assign{\Obj{0}}{M} \in \Domain{M}$, ${\prec} =
\Assign{<}{M} \subseteq M^2$, $* = \Assign{\prime}{M} \colon
\Domain{M} \to \Domain{M}$, and $\oplus = \Assign{+}{M}, \otimes =
\Assign{\times}{M} \colon \Domain{M}^2 \to \Domain{M}$. For each $x
\in \Domain{M}$, we write $x^*$ for the element of $\Domain{M}$
obtained from $x$ by application of~$*$.

Now, if $h$ were an isomorphism of $\Struct{N}$ and $\Struct{M}$,
there would be $n \in \Nat$ such that $h(n) = \Assign{c}{M}$.  So let
$s$ be any assignment in $\Struct{N}$ such that $s(x) = n$. Then
$\Sat{N}{\eq[\num{n}][x]}[s]$; by the proof of \olref[iso]{thm:isom},
also $\Sat{M}{\eq[\num{n}][x]}[h \circ s]$, so that $\Assign{c}{M} =
\mathbf{z}^{*\cdots *}$ (with $*$ iterated $n$ times). But this is impossible
since by assumption $\Sat{M}{\num{n} < c}$ and $\prec$ is
irreflexive. So $\Struct{M}$ is non-standard.
\end{proof}

\begin{prob}
A relation $R$ over a set $X$ is \emph{well-founded} if and only if
there are no infinite descending chains in~$R$, i.e., if there are no
$x_0$, $x_1$, $x_2$,~\dots in $X$ such that $\dots x_2Rx_1Rx_0$.  Assuming
Zermelo--Fraenkel set theory~$ZF$ is consistent, show that there are
non-well-founded models of $ZF$, i.e., models $\mathfrak{M}$ such that
$\dots x_2 \in x_1 \in x_0$.
\end{prob}

Since the non-standard model $\Struct{M}$ from \olref{thm:non-std} is
elementarily equivalent to the standard one, a number of properties of
$\Struct{M}$ can be derived. The rest of this section is devoted to
such a task, which will allow us to obtain a precise characterization
of \olup{enumerable} non-standard models of $\Theory{N}$.

\begin{enumerate}
\item No member of $\Domain{M}$ is $\prec$-less than itself: the sentence
  $\lforall[x][\lnot x < x]$ is true in $\Struct{N}$ and therefore in
  $\Struct{M}$.
\item By a similar reasoning we obtain that $\prec$ is a \emph{linear
    ordering} of $\Domain{M}$, i.e., a total, irreflexive, transitive relation
  on $\Domain{M}$. 
\item The element $\mathbf{z}$ is the $\prec$-least element of $\Domain{M}$.
\item Any member of $\Domain{M}$ is $\prec$-less than its $*$-successor and
  $x^*$ is the $\prec$-least member of $\Domain{M}$ greater than $x$.
\item $\Struct{M}$ contains an initial segment (of $\prec$) isomorphic
  to $\Nat$: $\mathbf{z}, \mathbf{z}^*, \mathbf{z}^{**}, \dots$, which
  we call the \emph{standard part} of $\Domain{M}$. Any other member
  of $\Domain{M}$ is \emph{non-standard}. There must be non-standard
  members of $\Domain{M}$, or else the function $h$ from the proof of
  \olref{thm:non-std} is an isomorphism.  We use $n, m, \dots$ as
  \olup{variable}s ranging on this standard part of $\Struct{M}$.
\item Every non-standard element is greater than any standard one;
  this is because for every $n \in \Nat$,
  \[
  \Sat{N}{\lforall[z][(\lnot(\eq[z][\Obj{0}] \lor
  \dots \lor \eq[z][\num{n}]) \lif \num{n} < z)]},
  \]
  so if $z \in \Domain{M}$ is different from all the standard
  elements, it must be \emph{greater} than all of them.
\item Any member of $\Domain{M}$ other than $\mathbf{z}$ is the
  $*$-successor of some unique element of $\Domain{M}$, denoted by
  $^*x$. If $x = y^*$ then both $x$ and $y$ are standard if one of
  them is (and both non-standard if one of them is).
\item Define an equivalence relation $\approx$ over $\Domain{M}$ by
  saying that $x \approx y$ if and only if for some \emph{standard}
  $n$, either $x \oplus n = y$ or $y \oplus n =x$. In other words, $x
  \approx y$ if and only if $x$ and $y$ are a finite distance
  apart. If $n$ and $m$ are standard then $n \approx m$. Define the
  \emph{block} of $x$ to be the equivalence class $[x] = \Setabs{y}{x
  \approx y}$.
\item Suppose that $x \prec y$ where $x \not\approx y$. Since
  $\Sat{N}{\lforall[x][\lforall[y][(x < y \lif (x' < y \lor x' =
      y))]]}$, either $x^* \prec y$ or $x^* = y$. The latter is
  impossible because it implies $x \approx y$, so $x \prec
  y$. Similarly, if $x \prec y$ and $x \not\approx y$, then $x \prec
  {^*y}$. Therefore if $x \prec y$ and $x \not\approx y$, then every
  $w \approx x$ is $\prec$-less than every $v \approx y$. Accordingly,
  each block $[x]$ forms a doubly infinite chain
  \[
  \cdots \prec  {^{**}x} \prec {^*}x \prec x \prec x^* \prec x^{**}
  \prec \cdots
  \]
  which is referred to as a $Z$-chain because it has the order type of
  the integers.
\item The $\prec$ ordering can be lifted up to the blocks: if $x \prec y$
  then the block of $x$ is less than the block of $y$. A block is
  \emph{non-standard} if it contains a non-standard element. The
  standard block is the least block. 
\item There is no least non-standard block: if $y$ is non-standard
  then there is a $x \prec y$ where $x$ is also non-standard and $x
  \not\approx y$. Proof: in the standard model $\Struct{N}$, every
  number is divisible by two, possibly with remainder one, i.e.,
  $\Sat{N}{\lforall[y][\lforall[x][(y = x + x \lor y = x + x +
        \Obj{0}')]]}$. By elementary equivalence, for every $y \in
  \Domain{M}$ there is $x \in \Domain{M}$ such that either $x \oplus x
  = y$ or $x \oplus x \oplus \mathbf{z}^*= y$. If $x$ were standard,
  then so would be $y$; so $x$ is non-standard. Furthermore, $x$ and
  $y$ belong to different blocks, i.e, $x \not\approx y$.  To see
  this, assume they did belong to the same block, i.e., $x \oplus n =
  y$ for some standard $n$. If $y = x \oplus x$, then $x \oplus n = x
  \oplus x$, whence $x = n$ by the cancellation law for addition
  (which holds in $\Struct{N}$ and therefore in $\Struct{M}$ as well),
  and $x$ would be standard after all. Similarly if $y = x \oplus x
  \oplus \mathbf{z}^*$.
\item By a similar argument, there is no greatest block. 
\item The ordering of the blocks is dense: if $[x]$ is less than $[y]$
  (where $x \not\approx y$), then there is a block $[z]$ distinct from
  both that is between them. Suppose $x \prec y$. As before, $x \oplus
  y$ is divisible by two (possibly with remainder) so there is a $u
  \in \Domain{M}$ such that either $x \oplus y = u \oplus u$ or $x
  \oplus y = u \oplus u \oplus \mathbf{z}^*$. The element $u$ is the
  average of $x$ and $y$, and so is between them. Assume $x \oplus y =
  u \oplus u$ (the other case being similar): if $u \approx x$ then
  for some standard $n$:
  \[
  x \oplus y = x \oplus n \oplus x \oplus n,
  \]
  so $y = x \oplus n \oplus n$ and we would have $x \approx y$,
  against assumption. We conclude that $u \not\approx x$. A similar
  argument gives $u \not\approx y$.
\end{enumerate}
The non-standard blocks are therefore ordered like the rationals: they
form \olup{enumerable} linear ordering without endpoints.  It follows
that for any two \olup{enumerable} non-standard models $\Struct{M}_1$ and
$\Struct{M_2}$ of true arithmetic, their reducts to the language
containing $<$ and $=$ only are isomorphic. Indeed, an isomorphism $h$
can be defined as follows: the standard parts of $\Struct{M_1}$ and
$\Struct{M_2}$ are isomorphic to the standard model $\Struct{N}$ and
hence to each other. The blocks making up the non-standard part are
themselves ordered like the rationals and therefore by
\olref[dlo]{thm:cantorQ} are isomorphic; an isomorphism of the blocks
can be extended to an isomorphism \emph{within} the blocks by matching
up arbitrary elements in each, and then taking the image of the
successor of $x$ in $\Struct{M_1}$ to be the successor of the image of
$x$ in $\Struct{M_2}$. Note that it does \emph{not} follow that
$\mathfrak{M}_1$ and $\mathfrak{M}_2$ are isomorphic in the full
language of arithmetic (indeed, isomorphism is always relative to a
signature), as there are non-isomorphic ways to define addition and
multiplication over $\Domain{M_1}$ and $\Domain{M_2}$. (This also
follows from a famous theorem due to Vaught that the number of
countable models of a complete theory cannot be~2.)

\begin{prob}
Show that there can be no greatest block in a non-standard model of
arithmetic.
\end{prob}

\begin{prob} 
Let $\Lang{L}$ be the first-order \olup{language} containing $<$ as its
only \olup{predicate} (besides $\eq$), and let $\Struct{N} = (\Nat,
<)$. All the finite or cofinite subsets of $\Struct{N}$ are
definable. Show that these are the \emph{only} definable subsets of
$\Struct{N}$.
 
(Hint: First, let $\Obj{prc}(x,y)$ be the
$\Lang{L}$-formula abbreviating ``$x$ is the immediate predecessor
of $y$:''
\[
x<y \land \lnot \lexists[z][(x<z \land z < y)].
\]
Now, to any definable subset of $\Struct{N}$ there corresponds a
formula $A(x)$ in $\Lang{L}$. For any such $A$,
consider the sentence $D$:
\[
\lexists[x][\lforall[y][\lforall[z][((x<y \land x<z \land \Obj{prc}(y,z)
  \land A(y)) \lif A(z))]]].
\]
Show that $\Sat{N}{D}$ if and only if the subset of
$\Struct{N}$ defined by $A$ is either finite or cofinite.

Now, let $\Struct{M}$ be a non-standard model elementarily
equivalent to $\Struct{N}$.  If $a \in \Domain{M}$ is
non-standard, let $b, c \in \Domain{M}$ be greater than $a$, and
let $b$ be the immediate predecessor of~$c$. Then there is an
automorphism $h$ of $\Domain{M}$ such that $h(b)=c$
(why?). Therefore, if $b$ satisfies $A$, so does $c$ (why?). It
follows that $D$ is true in $\Struct{M}$, and hence also in
$\Struct{N}$. But this implies that the subset of $\Struct{N}$
defined by $A$ is either finite or co-finite.
\end{prob}


% ==== OLP file: content/model-theory/basics/overspill.tex ====
% Part: first-order-logic
% Chapter: model-theory
% Section: overspill



\olfileid{mod}{bas}{ove}
\olsection{Overspill}

\begin{thm}
\ollabel{overspill} If a set $\Gamma$ of sentences has arbitrarily
large finite models, then it has an infinite model.
\end{thm}

\begin{proof}
Expand the language of $\Gamma$ by adding countably many new constants
$c_0$, $c_1$, \dots and consider the set $\Gamma \cup \{c_i \neq c_j :
i \neq j\}$. To say that $\Gamma$ has arbitrarily large finite models
means that for every $m >0$ there is $n\ge m$ such that $\Gamma$ has a
model of cardinality~$n$. This implies that $\Gamma \cup \{c_i \neq
c_j : i \neq j\}$ is finitely satisfiable. By compactness, $\Gamma
\cup \{c_i \neq c_j : i \neq j\}$ has a model $\Struct M$ whose domain
must be infinite, since it satisfies all inequalities $c_i \neq c_j$.
\end{proof}

\begin{prop}
\ollabel{inf-not-fo} 
There is no sentence $A$ of any first-order language that is true in
\olup{structure}~$\Struct M$ if and only if the domain $\Domain{M}$ of
the \olup{structure} is infinite.
\end{prop}

\begin{proof}
If there were such a $A$, its negation $\lnot A$ would be true in
all and only the finite \olup{structure}s, and it would therefore have
arbitrarily large finite models but it would lack an infinite model,
contradicting \olref{overspill}.
\end{proof}


% ==== OLP file: content/model-theory/basics/partial-iso.tex ====
% Part: first-order-logic
% Chapter: model-theory
% Section: partial-iso



\olfileid{mod}{bas}{pis}
\section{Partial Isomorphisms}

\begin{defn}
  Given two \olup{structure}s $\Struct{M}$ and $\Struct{N}$, a
  \emph{partial isomorphism} from $\Struct{M}$ to $\Struct{N}$ is a
  finite partial function $p$ taking arguments in $\Domain M$ and returning
  values in $\Domain N$, which satisfies the isomorphism conditions from
  \olref[iso]{defn:isomorphism} on its domain:
  \begin{enumerate}
  \item $p$ is \olup{injective};
  \item for every \olup{constant}~$c$: if $p(\Assign{c}{M})$ is defined,
    then $p(\Assign{c}{M}) = \Assign{c}{N}$;
  \item for every $n$-place \olup{predicate} $P$: if $a_1$, \dots, $a_n$
    are in the domain of $p$, then $\langle a_1, \dots, a_n\rangle \in
    \Assign P M$ if and only if $\langle p(a_1), \dots, p(a_n) \rangle
    \in \Assign P N$;
  \item for every $n$-place \olup{function} $f$: if $a_1$, \dots, $a_n$
    are in the domain of $p$, then $p(\Assign f M (a_1, \dots,a_n))
    = \Assign f N (p(a_1), \dots, p(a_n))$.
  \end{enumerate}
  That $p$ is finite means that $\dom{p}$ is finite.
\end{defn}

Notice that the empty function~$\emptyset$ is always a partial
isomorphism between any two \olup{structure}s.

\begin{defn}\ollabel{defn:partialisom}
  Two \olup{structure}s $\Struct{M}$ and $\Struct{N}$, are
  \emph{partially isomorphic}, written $\Struct{M} \iso[p]
  \Struct{N}$, if and only if there is a non-empty set $I$
  of partial isomorphisms between $\Struct{M}$ and $\Struct{N}$
  satisfying the \emph{back-and-forth} property:
  \begin{enumerate}
  \item (\emph{Forth}) For every $p \in I$ and $a \in \Domain M$
    there is $q \in I$ such that $p \subseteq q$ and $a$ is
    in the domain of $q$;
  \item (\emph{Back}) For every $p \in I$ and $b \in \Domain N$
    there is $q \in I$ such that $p \subseteq q$ and $b$ is
    in the range of $q$.
  \end{enumerate}
\end{defn}

\begin{thm}\ollabel{thm:p-isom1}
  If $\Struct{M} \iso[p] \Struct{N}$ and $\Struct{M}$ and
  $\Struct{N}$ are \olup{enumerable}, then $\Struct{M} \iso
  \Struct{N}$.
\end{thm}

\begin{proof}
  Since $\Struct{M}$ and $\Struct{N}$ are \olup{enumerable}, let $\Domain{M} =
  \{a_0, a_1, \ldots \}$ and $\Domain{N} = \{b_0, b_1, \ldots \}$. Starting
  with an arbitrary $p_0 \in I$, we define an increasing
  sequence of partial isomorphisms $p_0 \subseteq p_1 \subseteq p_2
  \subseteq \cdots$ as follows:
  \begin{enumerate}
  \item if $n+1$ is odd, say $n = 2r$, then using the Forth property
    find a $p_{n+1} \in I$ such that $p_n \subseteq p_{n+1}$
    and $a_r$ is in the domain of $p_{n+1}$;
  \item if $n+1$ is even, say $n+1 =2r$, then using the Back property
    find a $p_{n+1} \in I$ such that $p_n \subseteq p_{n+1}$
    and $b_r$ is in the range of $p_{n+1}$.
  \end{enumerate}
If we now put:
\[
p = \bigcup_{n\ge 0} p_n,
\]
we have that $p$ is a an isomorphism between $\Struct{M}$ and
$\Struct{N}$.
\end{proof}

\begin{prob}
  Show in detail that $p$ as defined in
  \olref[mod][bas][pis]{thm:p-isom1} is in fact an isomorphism.
\end{prob}

\begin{thm}\ollabel{thm:p-isom2}
  Suppose $\Struct{M}$ and $\Struct{N}$ are \olup{structure}s for a
  purely relational \olup{language} (\olup{language} containing only
  \olup{predicate}s, and no \olup{function}s or constants). Then if
  $\Struct{M} \iso[p] \Struct{N}$, also $\Struct{M} \elemequiv
  \Struct{N}$.
\end{thm}

\begin{proof}
  By induction on \olup{formula}s, one shows that if $a_1$, \dots, $a_n$ and
  $b_1$, \dots, $b_n$ are such that there is a partial isomorphism $p$
  mapping each $a_i$ to $b_i$ and $s_1(x_i) =a_i$ and $s_2(x_i) =b_i$
  (for $i =1$, \dots,~$n$), then $\Sat{M}{A}[s_1]$ if
  and only if $\Sat{N}{A}[s_2]$. The case for $n=0$
  gives $\Struct{M} \elemequiv \Struct{N}$.
\end{proof}

\begin{rem}
If \olup{function}s are present, the previous result is still true, but
one needs to consider the isomorphism induced by $p$ between the
sub\olup{structure} of $\Struct{M}$ generated by $a_1$, \dots, $a_n$ and the
sub\olup{structure} of $\Struct{N}$ generated by $b_1$, \dots, $b_n$.
\end{rem}

The previous result can be ``broken down'' into stages by establishing a
connection between the number of nested quantifiers in \olup{formula} and
how many times the relevant partial isomorphisms can be extended.

\begin{defn}
  For any \olup{formula}~$A$, the \emph{quantifier rank} of $A$, denoted
  by $\QuantRank{A} \in \Nat$, is recursively defined as
  the highest number of nested quantifiers in $A$.  Two
  \olup{structure}s $\Struct{M}$ and $\Struct{N}$ are \emph{$n$-equivalent},
  written $\Struct{M} \elemequiv[n] \Struct{N}$, if they agree on all
  sentences of quantifier rank less than or equal to~$n$.
\end{defn}

\begin{prop}\ollabel{prop:qr-finite}
  Let $\Lang{L}$ be a finite purely relational \olup{language}, i.e., a
  \olup{language} containing finitely many \olup{predicate}s and \olup{constant}s,
  and no \olup{function}s. Then for each $n \in \Nat$ there are
  only finitely many first-order \olup{sentence}s in the \olup{language}
  $\Lang{L}$ that have quantifier rank no greater than $n$, up to
  logical equivalence.
\end{prop}

\begin{proof}
  By induction on $n$.
\end{proof}

\begin{defn}
  Given \olup{structure}~$\Struct{M}$, let $\Domain M^{<\omega}$ be the set of
  all finite sequences over $\Domain{M}$. We use $\mathbf{a},
  \mathbf{b}, \mathbf{c}, \ldots$ to range over finite sequences of
  elements. If $\mathbf{a} \in \Domain{M}^{<\omega}$ and $a \in \Domain{M}$, then
  $\mathbf{a}a$ represents the \emph{concatenation} of $\mathbf{a}$ with $a$.
\end{defn}

\begin{defn}
  Given \olup{structure}s $\Struct{M}$ and $\Struct{N}$, we define
  relations $I_n \subseteq \Domain M^{<\omega} \times \Domain N^{<\omega}$ between
  sequences of equal length, by recursion on $n$ as follows:
   \begin{enumerate}
   \item $I_0(\mathbf{a},\mathbf{b})$ if and only if $\mathbf{a}$ and
     $\mathbf{b}$ satisfy the same atomic \olup{formula}s in  $\Struct{M}$
     and  $\Struct{N}$; i.e., if $s_1(x_i) = a_i$ and $s_2(x_i) =
     b_i$ and $A$ is atomic with all \olup{variable}s among
     $x_1$, \dots,~$x_n$, then $\Sat{M}{A}[s_1]$ if and
     only if~$\Sat{N}{A}[s_2]$.
   \item $I_{n+1} (\mathbf{a},\mathbf{b})$ if and only if for every
     $a\in \Domain M$ there is a $b\in \Domain N$ such that $I_n
     (\mathbf{a}a,\mathbf{b}b)$, and vice-versa.
   \end{enumerate}
\end{defn}


\begin{defn}
  Write $\Struct{M} \approx_n \Struct{N}$ if
  $I_n(\emptyseq,\emptyseq)$ holds of $\Struct{M}$ and
  $\Struct{N}$ (where $\emptyseq$ is the empty sequence).
\end{defn}

\begin{thm}\ollabel{thm:b-n-f}
  Let $\Lang{L}$ be a purely relational \olup{language}. Then $I_n
  (\mathbf{a},\mathbf{b})$ implies that for every $A$ such that
  $\QuantRank{A} \le n$, we have $\Sat{M}{A}[\mathbf{a}]$ if and
  only if $\Sat{N}{A}[\mathbf{b}]$ (where again $\mathbf{a}$
  satisfies $A$ if any $s$ such that $s(x_i) = a_i$ satisfies
  $A$). Moreover, if $\Lang{L}$ is finite, the converse also holds.
\end{thm}

\begin{proof}
  The proof that $I_n(\mathbf{a},\mathbf{b})$ implies that
  $\mathbf{a}$ and $\mathbf{b}$ satisfy the same \olup{formula}s of
  quantifier rank no greater than $n$ is by an easy induction on
  $A$. For the converse we proceed by induction on $n$, using
  \olref{prop:qr-finite}, which ensures that for each $n$
  there are at most finitely many non-equivalent \olup{formula}s of that
  quantifier rank.

  For $n=0$ the hypothesis that $\mathbf{a}$ and $\mathbf{b}$ satisfy
  the same quantifier-free \olup{formula}s gives that they satisfy the same
  atomic ones, so that $I_0(\mathbf{a},\mathbf{b})$.

  For the $n+1$ case, suppose that $\mathbf{a}$ and $\mathbf{b}$
  satisfy the same \olup{formula}s of quantifier rank no greater than
  $n+1$; in order to show that $I_{n+1}(\mathbf{a},\mathbf{b})$
  suffices to show that for each $a \in \Domain M$ there is a $b \in
  \Domain N$ such that $I_n(\mathbf{a}a,\mathbf{b}b)$, and by the
  inductive hypothesis again suffices to show that for each $a \in
  \Domain M$ there is a $b \in \Domain N$ such that $\mathbf{a}a$ and
  $\mathbf{b}b$ satisfy the same \olup{formula}s of quantifier rank no
  greater than $n$.

  Given $a \in \Domain M$, let $T^a_n$ be set of \olup{formula}s
  $B(x,\mathbf{y})$ of rank no greater than $n$ satisfied by
  $\mathbf{a}a$ in $\Struct{M}$; $T^a_n$ is finite, so we can
  assume it is a single first-order \olup{formula}. It follows that
  $\mathbf{a}$ satisfies $\lexists[x][T^a_n(x,\mathbf{y})]$, which
  has quantifier rank no greater than $n+1$. By hypothesis
  $\mathbf{b}$ satisfies the same \olup{formula} in $\Struct{N}$, so that
  there is a $b \in \Domain N$ such that $\mathbf{b}b$ satisfies
  $T^a_n$; in particular, $\mathbf{b}b$ satisfies the same
  \olup{formula}s of quantifier rank no greater than $n$ as
  $\mathbf{a}a$. Similarly one shows that for every $b \in \Domain N$
  there is $a\in \Domain M$ such that $\mathbf{a}a$ and $\mathbf{b}b$
  satisfy the same \olup{formula}s of quantifier rank no greater than $n$,
  which completes the proof.
\end{proof}

\begin{cor}\ollabel{cor:b-n-f}
  If $\Struct{M}$ and $\Struct{N}$ are purely relational \olup{structure}s
  in a finite \olup{language}, then $\Struct{M} \approx_n\Struct{N}$ if and
  only if $\Struct{M} \elemequiv[n] \Struct{N}$. In particular
  $\Struct{M} \elemequiv \Struct{N}$ if and only if for each $n$,
  $\Struct{M} \approx_n \Struct{N}$ .
\end{cor}


% ==== OLP file: content/model-theory/basics/reducts-and-expansions.tex ====
% Part: first-order-logic
% Chapter: model-theory
% Section: reducts-and-expansions



\olfileid{mod}{bas}{red}
\section{Reducts and Expansions}

Often it is useful or necessary to compare languages which have
symbols in common, as well as \olup{structure}s for these languages.  The
most common case is when all the symbols in \olup{language}~$\Lang{L}$
are also part of \olup{language}~$\Lang{L'}$, i.e., $\Lang{L} \subseteq
\Lang{L'}$. An $\Lang{L}$-\olup{structure}~$\Struct{M}$ can then always
be expanded to an $\Lang{L'}$-\olup{structure} by adding interpretations
of the additional symbols while leaving the interpretations of the
common symbols the same.  On the other hand, from an
$\Lang{L'}$-structure~$\Struct{M'}$ we can obtain an
$\Lang{L}$-structure simply by ``forgetting'' the interpretations of
the symbols that do not occur in~$\Lang{L}$.

\begin{defn}
\ollabel{defn:reduct}
Suppose $\Lang L \subseteq \Lang L'$, $\Struct M$ is an
$\Lang L$-\olup{structure} and $\Struct M'$ is an $\Lang L'$-\olup{structure}.
$\Struct M$ is the \emph{reduct} of $\Struct M'$ to $\Lang L$, and
$\Struct M'$ is an \emph{expansion} of $\Struct M$ to $\Lang L'$ iff
\begin{enumerate}
\item $\Domain{M} = \Domain{M'}$
\item For every \olup{constant}~$c \in \Lang L$, $\Assign{c}{M} =
  \Assign{c}{M'}$.
\item For every \olup{function}~$f \in \Lang L$, $\Assign{f}{M} =
  \Assign{f}{M'}$.
\item For every \olup{predicate}~$P \in \Lang L$, $\Assign{P}{M} =
  \Assign{P}{M'}$.
\end{enumerate}
\end{defn}

\begin{prop}
\ollabel{prop:reduct}
If an $\Lang{L}$-\olup{structure}~$\Struct{M}$ is a reduct of an
$\Lang{L'}$-\olup{structure}~$\Struct{M'}$, then for all
$\Lang{L}$-\olup{sentence}s~$A$,
\[
\Sat{M}{A} \text{ iff } \Sat{M'}{A}.
\]
\end{prop}

\begin{proof}
  Exercise.
\end{proof}

\begin{prob}
Prove \olref[mod][bas][red]{prop:reduct}.
\end{prob}

\begin{defn}
When we have an $\Lang{L}$-structure $\Struct{M}$, and $\Lang{L'} =
\Lang{L} \cup \{P\}$ is the expansion of $\Lang{L}$ obtained by adding
a single $n$-place \olup{predicate}~$P$, and $R \subseteq \Domain{M}^n$
is an $n$-place relation, then we write $\Expan{M}{R}$ for the
expansion~$\Struct{M'}$ of~$\Struct{M}$ with $\Assign{P}{M'} = R$.
\end{defn}


% ==== OLP file: content/model-theory/basics/substructures.tex ====
% Part: first-order-logic
% Chapter: model-theory
% Section: substructures



\olfileid{mod}{bas}{sub}
\olsection{Sub\printtoken{p}{structure}}

The \olup{domain} of \olup{structure}~$\Struct{M}$ may be a subset of
another~$\Struct{M'}$.  But we should obviously only consider
$\Struct{M}$ a ``part'' of $\Struct{M'}$ if not only $\Domain{M}
\subseteq \Domain{M'}$, but $\Struct{M}$ and $\Struct{M'}$ ``agree''
in how they interpret the symbols of the language at least on the
shared part~$\Domain{M}$.

\begin{defn}
\ollabel{defn:substructure}
Given \olup{structure}s $\Struct M$ and $\Struct M'$ for the same
language~$\Lang L$, we say that $\Struct M$ is a \emph{sub\olup{structure}}
of $\Struct M'$, and $\Struct M'$ an \emph{extension} of $\Struct M$,
written $\Struct M \substruct \Struct M'$, iff
\begin{enumerate}
\item $\Domain{M} \subseteq \Domain{M'}$,
\item For each constant $c \in \Lang L$, $\Assign{c}{M} =
    \Assign{c}{M'}$;
\item For each $n$-place \olup{function} $f \in \Lang L$
  $\Assign{f}{M}(a_1, \dots, a_n) = \Assign{f}{M'}(a_1, \dots, a_n)$
  for all $a_1$, \dots, $a_n \in \Domain{M}$.
\item For each $n$-place \olup{predicate} $R \in \Lang L$, $\langle
  a_1, \dots, a_n\rangle \in \Assign{R}{M}$ iff $\langle a_1, \dots,
  a_n\rangle \in \Assign{R}{M'}$ for all $a_1$, \dots, $a_n \in
  \Domain{M}$.
\end{enumerate}
\end{defn}

\begin{rem}
\ollabel{rem:substructure}
If the language contains no constant or \olup{function}s, then any $N
\subseteq \Domain{M}$ determines a sub\olup{structure}~$\Struct{N}$ of
$\Struct M$ with \olup{domain}~$\Domain{N} = N$ by putting $\Assign{R}{N} =
\Assign{R}{M} \cap N^n$.
\end{rem}

% prove something about this? Examples?


% ==== OLP file: content/model-theory/basics/theory-of-m.tex ====
% Part: first-order-logic
% Chapter: model-theory
% Section: theory-of-m



\olfileid{mod}{bas}{thm}

\section{The Theory of a \printtoken{S}{structure}}

Every \olup{structure}~$\Struct{M}$ makes some \olup{sentence}s true, and
some false. The set of all the \olup{sentence}s it makes true is called
its \emph{theory}. That set is in fact a theory, since anything it
entails must be true in all its models, including~$\Struct{M}$.

\begin{defn}
  Given \olup{structure}~$\Struct M$, the \emph{theory} of
  $\Struct{M}$ is the set $\Theory{M}$ of \olup{sentence}s
  that are true in $\Struct{M}$, i.e., $\Theory{M} =
  \Setabs{A}{\Sat{M}{A}}$.
\end{defn}

We also use the term ``theory'' informally to refer to sets
of \olup{sentence}s having an intended interpretation, whether deductively
closed or not.

\begin{prop}
For any $\Struct{M}$, $\Theory{M}$ is complete.
\end{prop}

\begin{proof}
For any \olup{sentence}~$A$ either $\Sat{M}{A}$ or $\Sat{M}{\lnot A}$,
so either $A \in \Theory{M}$ or $\lnot A \in \Theory{M}$.
\end{proof}

\begin{prop}\ollabel{prop:equiv}
  If $\Struct{N} \models A$ for every $A \in \Theory{M}$, then
  $\Struct{M} \elemequiv \Struct{N}$.
\end{prop}

\begin{proof}
Since $\Sat{N}{A}$ for all $A \in \Theory{M}$, $\Theory{M} \subseteq
\Theory{N}$. If $\Sat{N}{A}$, then $\Sat/{N}{\lnot A}$, so $\lnot A
\notin \Theory{M}$. Since $\Theory{M}$ is complete, $A \in
\Theory{M}$. So, $\Theory{N} \subseteq \Theory{M}$, and we have
$\Struct{M} \elemequiv \Struct{N}$.
\end{proof}

\begin{rem}\ollabel{remark:R}
  Consider $\Struct{R} = \langle\Real, <\rangle$, the \olup{structure}
  whose domain is the set $\Real$ of the real numbers, in the \olup{language}
  comprising only a 2-place \olup{predicate} interpreted as the $<$
  relation over the reals. Clearly $\Struct{R}$ is \olup{nonenumerable};
  however, since $\Theory{R}$ is obviously consistent, by the
  L\"owenheim--Skolem theorem it has \olup{enumerable} model, say
  $\Struct{S}$, and by \olref{prop:equiv}, $\Struct{R}
  \equiv \Struct{S}$. Moreover, since $\Struct{R}$ and $\Struct{S}$
  are not isomorphic, this shows that the converse of
  \olref[iso]{thm:isom} fails in general.
\end{rem}


\end{document}
